\documentclass[11pt,a4paper]{article}

% ============================================================
% PACKAGES
% ============================================================
\usepackage[utf8]{inputenc}
\usepackage[T1]{fontenc}
\usepackage[english]{babel}
\usepackage{amsmath,amssymb,amsthm,mathtools}
\usepackage{mathrsfs}
\usepackage{tikz-cd}
\usepackage{enumitem}
\usepackage{hyperref}
\usepackage{geometry}
\usepackage{microtype}
\usepackage{booktabs}

\geometry{margin=1in}

% ============================================================
% THEOREM ENVIRONMENTS (inherited)
% ============================================================
\theoremstyle{definition}
\newtheorem{definition}{Definition}[section]
\newtheorem{axiom}[definition]{Axiom}
\newtheorem{example}[definition]{Example}
\newtheorem{remark}[definition]{Remark}
\newtheorem{convention}[definition]{Convention}
\newtheorem{construction}[definition]{Construction}
\newtheorem{problem}[definition]{Problem}

\theoremstyle{plain}
\newtheorem{proposition}[definition]{Proposition}
\newtheorem{theorem}[definition]{Theorem}
\newtheorem{lemma}[definition]{Lemma}
\newtheorem{corollary}[definition]{Corollary}

% ============================================================
% MACROS (inherited from Papers I / framework)
% ============================================================
\newcommand{\Lam}{\Lambda}
\newcommand{\Sp}{\mathrm{Sp}}
\newcommand{\Ge}{\mathrm{Ge}}
\newcommand{\To}{\mathrm{To}}
\newcommand{\Al}{\mathrm{Al}}
\newcommand{\Nu}{\mathrm{Nu}}
\newcommand{\Fo}{\mathrm{Fo}}
\newcommand{\Cl}{\mathrm{Cl}}
\newcommand{\Corr}{\mathrm{Corr}}
\newcommand{\An}{\mathrm{An}}
\newcommand{\Ob}{\mathfrak{O}}
\newcommand{\Ocal}{\mathcal{O}}
\newcommand{\Rcal}{\mathcal{R}}
\newcommand{\Ccal}{\mathcal{C}}
\newcommand{\Xcal}{\mathcal{X}}
\newcommand{\Ycal}{\mathcal{Y}}
\newcommand{\Pcal}{\mathcal{P}}
\newcommand{\Mfrak}{\mathfrak{M}}
\newcommand{\Rfrak}{\mathfrak{R}}
\newcommand{\Phimap}{\Phi}
\newcommand{\simadm}{\sim}
\newcommand{\dd}{\,\mathrm{d}}
\newcommand{\op}[1]{\operatorname{#1}}
\newcommand{\spec}{\sigma}
\newcommand{\Norm}[1]{\left\lVert #1 \right\rVert}
\newcommand{\Abs}[1]{\left\lvert #1 \right\rvert}
\newcommand{\R}{\mathbb{R}}
\newcommand{\C}{\mathbb{C}}
\newcommand{\N}{\mathbb{N}}
\newcommand{\Z}{\mathbb{Z}}

% ============================================================
% TITLE
% ============================================================
\title{%
  \textbf{On the Limits of Structural Constraining:}\\[4pt]
  \large Separation, Coherence, and the Emergence of New Representations
}
\author{}
\date{}

\begin{document}
\maketitle

\begin{abstract}
\noindent
The companion paper \cite{paperI} develops a
multirepresentational methodology in which a mathematical model is
progressively constrained by successive layers---spectral,
geometric, topological, algebraic, numerical, analytic---until the
admissible class reduces to the model itself. The present paper
addresses a question left open there: \emph{how far can the
constraining procedure go?} We formalise the closure of the
methodology as the injectivity of a representation map $\Phimap$
modulo an admissible equivalence relation, and we classify the
possible failures into three types: representational rupture
(recoverable by emergence), informational rupture (recoverable by
adding independent representations), and fundamental rupture
(recoverable only by introducing genuinely new structure). We
formalise the cross-coherence of layers as a subspace $\Rfrak_{\rm
coh} \subset \prod_i \Rcal_i$ rather than a product, and we give a
dynamic account of representation expansion in which the residue of
the constraining chain generates pressure for the introduction of
the next representation. We close with the open problem of
\emph{coherent indistinguishability}: does there exist a pair of
non-isomorphic mathematical structures indistinguishable under all
coherent representations?
\end{abstract}

\tableofcontents

% ------------------------------------------------------------
\section{Introduction}
\label{sec:intro}
% ------------------------------------------------------------

\subsection{Context}

The companion paper \cite{paperI} introduces a multirepresentational
methodology in which a mathematical model $P$ is constrained
successively by spectral, geometric, topological, algebraic,
numerical, analytic, and formal layers. Each layer $\lambda$ defines
a constrained model space $\Ccal_\lambda$, and the methodology's
output is the intersection
\[
  \Ccal_\ast = \bigcap_{\lambda \in \Lam} \Ccal_\lambda .
\]

The companion paper establishes three structural principles:
\emph{constraining}, \emph{spectral emergence}, and
\emph{discrepancy classification}. It also gives four worked
examples (Dirichlet Laplacian, path-vs-cycle bifurcation, erroneous
diffusion, nonlinear $p$-Laplacian).

\subsection{The question left open}

The companion paper does not address a foundational question: 
\emph{is the constraining procedure guaranteed to determine the
model?} That is, does the intersection $\Ccal_\ast$ collapse to a
singleton (or to a single equivalence class), or can it remain
larger even after all layers have been applied?

The present paper answers this question by:
\begin{enumerate}[label=(\roman*)]
  \item formalising the closure criterion as the injectivity of a
        representation map (\S\ref{sec:separation});
  \item classifying the possible failures into three types
        (\S\ref{sec:rupture});
  \item formalising cross-coherence of layers
        (\S\ref{sec:coherence});
  \item giving a dynamic account of representation expansion
        (\S\ref{sec:expansion});
  \item stating the open problem of coherent indistinguishability
        (\S\ref{sec:open}).
\end{enumerate}

\subsection{Contributions}

\begin{enumerate}[label=(\roman*)]
  \item \textbf{Separation criterion} (\S\ref{sec:separation}): the
        methodology closes if and only if $\Phimap$ is injective
        modulo the admissible equivalence $\simadm$.
  \item \textbf{Three-type rupture classification}
        (\S\ref{sec:rupture}): a classification of failure modes
        with distinct diagnostic and remedial signatures.
  \item \textbf{Cross-coherence} (\S\ref{sec:coherence}): the
        admissible representation space is a subspace of the
        product, not the product itself.
  \item \textbf{Dynamic expansion} (\S\ref{sec:expansion}): the
        residue of the chain at level $n$ generates pressure for
        the introduction of the next representation.
  \item \textbf{Open problem} (\S\ref{sec:open}): the existence of
        a pair of non-isomorphic structures indistinguishable under
        all coherent representations.
\end{enumerate}

% ------------------------------------------------------------
\section{The Separation Criterion}
\label{sec:separation}
% ------------------------------------------------------------

\subsection{The Representation Map}

\begin{definition}[Representation family]\label{def:rep-family}
Let $\Xcal$ be a class of mathematical structures and let
\[
  \Rfrak = \{ R_i : \Xcal \to \Rcal_i \}_{i \in I}
\]
be a family of \emph{representations}, where each $\Rcal_i$ is a
codomain (a category, a set, a space, or a formal theory). The
\emph{representation map} associated with $\Rfrak$ is
\[
  \Phimap_{\Rfrak} : \Xcal \longrightarrow \prod_{i \in I} \Rcal_i,
  \qquad
  \Phimap_{\Rfrak}(X) = (R_i(X))_{i \in I}.
\]
\end{definition}

\begin{convention}
We write $\Phimap$ for $\Phimap_{\Rfrak}$ when $\Rfrak$ is clear
from context.
\end{convention}

\subsection{The Admissible Equivalence}

\begin{definition}[Admissible equivalence]\label{def:adm-equiv}
An \emph{admissible equivalence} on $\Xcal$ is an equivalence
relation $\simadm$ such that $X \simadm Y$ iff $X$ and $Y$ represent
the same mathematical structure modulo change of coordinates,
gauge, isomorphism, diffeomorphism, or any other transformation
declared to be structure-preserving.
\end{definition}

\begin{example}
For Riemannian manifolds, $\simadm$ is isometry. For operators on a
fixed Hilbert space, $\simadm$ is unitary equivalence. For groups,
$\simadm$ is isomorphism.
\end{example}

\subsection{The Closure Criterion}

\begin{definition}[Separation]\label{def:separation}
The family $\Rfrak$ \emph{separates} $\Xcal$ modulo $\simadm$ if
\[
  \Phimap(X) = \Phimap(Y)
  \;\Longrightarrow\;
  X \simadm Y .
\]
Equivalently, the induced map
\[
  \overline{\Phimap} : \Xcal / \simadm \longrightarrow \prod_i \Rcal_i
\]
is injective.
\end{definition}

\begin{theorem}[Closure criterion]\label{thm:closure}
The multirepresentational methodology of \cite{paperI} closes on a
class $\Xcal$ if and only if the representation family
$\Rfrak = \{R_\lambda\}_{\lambda \in \Lam}$ separates $\Xcal$
modulo $\simadm$.
\end{theorem}

\begin{proof}
$(\Rightarrow)$ Suppose the methodology closes: the constrained
model space $\Ccal_\ast$ reduces to a single equivalence class.
Then any two models $X, Y$ with the same representation data must
belong to the same constrained space, hence to the same
$\simadm$-class. Therefore $\Phimap(X) = \Phimap(Y)$ implies
$X \simadm Y$.

$(\Leftarrow)$ Suppose $\Phimap$ separates $\Xcal$ modulo
$\simadm$. Let $X \in \Ccal_\ast$. Every $Y \in \Ccal_\ast$ has
$\Phimap(Y) = \Phimap(X)$ by definition of $\Ccal_\ast$ (each layer
sees the same data). Separation then gives $Y \simadm X$, so
$\Ccal_\ast$ is a single equivalence class.
\end{proof}

\begin{remark}
The closure criterion is a statement about the \emph{family}
$\Rfrak$, not about any individual representation. A single layer
$R_i$ may fail to be injective; only the joint map $\Phimap$
matters.
\end{remark}

\subsection{The Diagonal Case}

\begin{proposition}[Failure of separation]\label{prop:failure-sep}
$\Rfrak$ fails to separate $\Xcal$ modulo $\simadm$ if and only if
there exist $X \not\simadm Y$ with
\[
  R_i(X) = R_i(Y) \qquad \forall i \in I .
\]
\end{proposition}

\begin{proof}
Immediate from Definition~\ref{def:separation}.
\end{proof}

\begin{remark}
A failure of separation on a \emph{single} representation
$R_i$---e.g.\ iso-spectral non-isometric manifolds, or two
nonlinear equations with the same linearisation---is \emph{not} a
failure of the methodology. It is a signal that $R_i$ alone is
insufficient, which is exactly what the methodology anticipates.
The genuine failure is the one stated in
Proposition~\ref{prop:failure-sep}: indistinguishability under
\emph{all} representations simultaneously.
\end{remark}

% ------------------------------------------------------------
\section{Three Types of Rupture}
\label{sec:rupture}
% ------------------------------------------------------------

The methodology may fail in three structurally distinct ways. The
distinction is diagnostic: each type has a different signature and
a different remedy.

\subsection{Type I: Representational Rupture}

\begin{definition}[Representational rupture]\label{def:rupture-I}
A layer $\lambda$ exhibits a \emph{representational rupture} if the
representation $R_\lambda$ does not exist in the initially sought
form.
\end{definition}

\begin{example}[Nonlinear spectral layer]
For a nonlinear evolution equation $u_t = \Ncal(u)$, no global
spectral layer exists in the linear sense. However, Fr\'echet
linearisation at an equilibrium $u_\ast$ produces an emergent
spectral layer
\[
  R_\Sp^{\rm emg} = \spec(D\Ncal[u_\ast]).
\]
The rupture is \emph{Type~I}: the sought representation is
constructed, not assumed.
\end{example}

\begin{proposition}[Remedy for Type I]\label{prop:remedy-I}
A representational rupture is remedied by construction: the missing
representation is produced from the existing layers via an
appropriate emergent construction (linearisation, completion,
localisation, or regularisation).
\end{proposition}

\begin{proof}
By definition of Type~I, the rupture concerns the \emph{existence}
of the representation. If a construction produces it, the rupture
is resolved. The framework paper's emergent spectral layer is the
canonical example.
\end{proof}

\subsection{Type II: Informational Rupture}

\begin{definition}[Informational rupture]\label{def:rupture-II}
A layer $\lambda$ exhibits an \emph{informational rupture} if
$R_\lambda$ exists but is not injective on $\Xcal$ modulo $\simadm$:
\[
  \exists X \not\simadm Y, \quad R_\lambda(X) = R_\lambda(Y).
\]
\end{definition}

\begin{example}[Isospectrality]
There exist Riemannian manifolds $M \not\cong N$ with
$\spec(\Delta_M) = \spec(\Delta_N)$ (Gordon--Webb--Wolpert). The
spectral layer is informative but not injective.
\end{example}

\begin{example}[Linearisation]
The equations
\[
  u_t = u_{xx} - u^3, \qquad u_t = u_{xx} - 2 u^3
\]
have the same Fr\'echet linearisation at $u = 0$, hence the same
emergent spectrum, but are distinct nonlinear evolution equations.
\end{example}

\begin{proposition}[Remedy for Type II]\label{prop:remedy-II}
An informational rupture is remedied by adding an independent
representation $R_\mu$ such that the pair $(R_\lambda, R_\mu)$ is
injective modulo $\simadm$ on $\Xcal$.
\end{proposition}

\begin{proof}
If $R_\lambda$ fails injectivity on a pair $(X, Y)$, we require
$R_\mu(X) \neq R_\mu(Y)$ for some additional representation
$R_\mu$. The combined map $X \mapsto (R_\lambda(X), R_\mu(X))$ is
then injective on that pair, and by iteration, on $\Xcal$.
\end{proof}

\begin{example}[Remedy for the linearisation example]
The variational layer distinguishes the two equations above: the
functionals
\[
  E_1[u] = \int \left( \tfrac{1}{2} u_x^2 + \tfrac{1}{4} u^4 \right),
  \quad
  E_2[u] = \int \left( \tfrac{1}{2} u_x^2 + \tfrac{1}{2} u^4 \right)
\]
differ. The pair (spectral, variational) separates them.
\end{example}

\subsection{Type III: Fundamental Rupture}

\begin{definition}[Fundamental rupture]\label{def:rupture-III}
A family $\Rfrak$ exhibits a \emph{fundamental rupture} if
\[
  \exists X \not\simadm Y : \quad R_i(X) = R_i(Y) \quad \forall i \in I,
\]
and no admissible extension of $\Rfrak$ within the declared
representation class separates $X$ from $Y$.
\end{definition}

\begin{remark}
Type~III is the genuine failure of the methodology. Type~I is
surmountable by construction; Type~II is surmountable by addition;
Type~III requires a genuinely new kind of representation, outside
the declared class.
\end{remark}

\subsection{Classification Table}

\begin{center}
\begin{tabular}{c|c|c|c}
  \toprule
  Type & Signature & Remedy & Status \\
  \midrule
  I & $R_\lambda$ does not exist & construction
    & surmountable \\
  II & $R_\lambda$ not injective & add $R_\mu$
    & surmountable \\
  III & $\Phimap$ not injective, class saturated & new class
    & open \\
  \bottomrule
\end{tabular}
\end{center}

\begin{proposition}[Hierarchy of ruptures]\label{prop:hierarchy}
Type~I rupture may give rise to Type~II (once the representation
exists, it may fail injectivity). Type~II rupture may give rise to
Type~III (if no admissible extension separates). The three types
are strictly nested:
\[
  \text{Type I} \;\subsetneq\; \text{Type II} \;\subsetneq\;
  \text{Type III}.
\]
\end{proposition}

\begin{proof}
Immediate from the definitions.
\end{proof}

% ------------------------------------------------------------
\section{Cross-Coherence Constraints}
\label{sec:coherence}
% ------------------------------------------------------------

\subsection{The Naive Product is Wrong}

A naive view treats the representations as independent:
\[
  \Phimap(X) \in \prod_{i \in I} \Rcal_i
  \qquad \text{(naive)}.
\]
But the representations of a mathematical object are not
independent. Spectral data constrain geometric data, geometric data
constrain topological data, topological data constrain algebraic
data, and so on. The admissible space is a \emph{subspace}, not the
full product.

\begin{definition}[Cross-coherence]\label{def:cross-coherence}
A \emph{coherence relation} on the family $\{R_i\}_{i \in I}$ is a
subset
\[
  \Rfrak_{\rm coh} \;\subset\; \prod_{i \in I} \Rcal_i
\]
defined by compatibility conditions of the form
\[
  C_{ij} \subset \Rcal_i \times \Rcal_j,
  \qquad (i, j) \in I \times I,
\]
such that
\[
  \Rfrak_{\rm coh}
  =
  \bigl\{ (r_i)_{i \in I} : (r_i, r_j) \in C_{ij} \;\;
  \forall i, j \bigr\}.
\]
\end{definition}

\begin{example}[Spectral--geometric coherence]
The isospectrality of a manifold determines a class of admissible
metrics. The pair (spectrum, metric) must lie in the coherence
subset determined by the inverse spectral problem:
\[
  C_{\Sp\Ge} = \{(\spec(\Delta_g), g) : g \text{ Riemannian}, \;
  \spec(\Delta_g) = \spec\}.
\]
\end{example}

\subsection{Constrained Representation}

\begin{definition}[Coherent representation map]\label{def:coherent-map}
The \emph{coherent representation map} is
\[
  \Phimap_{\rm coh} : \Xcal \longrightarrow \Rfrak_{\rm coh},
  \qquad
  \Phimap_{\rm coh}(X) = \Phimap(X) \text{ if } \Phimap(X) \in \Rfrak_{\rm coh}.
\]
\end{definition}

\begin{theorem}[Coherence strengthens separation]\label{thm:coherence-strengthens}
If $\Phimap$ fails to separate $\Xcal$ modulo $\simadm$ but
$\Phimap_{\rm coh}$ separates, then the constraining has been
completed by the coherence relations.
\end{theorem}

\begin{proof}
Suppose $\Phimap(X) = \Phimap(Y)$ with $X \not\simadm Y$. If
$\Phimap_{\rm coh}$ separates, then at least one component of the
tuple $\Phimap(X)$ violates a coherence constraint that
$\Phimap(Y)$ satisfies, so $\Phimap_{\rm coh}(X) \neq
\Phimap_{\rm coh}(Y)$.
\end{proof}

\begin{remark}
This is the formal content of the intuition that ``layers
\emph{force} each other.'' The coherence relations are what turn a
collection of invariants into a mutually constraining system.
\end{remark}

\subsection{Coherence as a Subspace}

The admissible space is
\[
  \Rfrak_{\rm coh}
  =
  \bigcap_{(i,j)} \pi_{ij}^{-1}(C_{ij}),
\]
where $\pi_{ij} : \prod_k \Rcal_k \to \Rcal_i \times \Rcal_j$ is
the projection. The coherence structure is therefore a
\emph{fibered} constraint system.

\begin{proposition}[Coherence is not a product]\label{prop:not-product}
In general, $\Rfrak_{\rm coh} \neq \prod_{i} \Rcal_i$.
\end{proposition}

\begin{proof}
The compatibility conditions $C_{ij}$ are nontrivial proper subsets
of $\Rcal_i \times \Rcal_j$ in any non-degenerate application.
\end{proof}

% ------------------------------------------------------------
\section{Dynamic Representation Expansion}
\label{sec:expansion}
% ------------------------------------------------------------

\subsection{The Residue}

\begin{definition}[Residue of the constraining chain]\label{def:residue}
Given a representation family $\Rfrak_n = \{R_1, \ldots, R_n\}$ and
observed data $(r_1, \ldots, r_n)$, the \emph{residue} is
\[
  \Ccal_\infty^{(n)}
  =
  \bigl\{ X \in \Xcal :
    \Phimap_{\Rfrak_n}(X) = (r_1, \ldots, r_n) \bigr\}
  \big/ \simadm .
\]
\end{definition}

The residue measures the information still undetermined by the
current representation family.

\subsection{Residue-Driven Expansion}

\begin{axiom}[Expansion axiom]\label{ax:expansion}
If the residue $\Ccal_\infty^{(n)}$ contains more than one
equivalence class, the methodology is \emph{not closed}. The next
representation $R_{n+1}$ is chosen so as to reduce the residue:
\[
  \Rfrak_{n+1} = \Rfrak_n \cup \{R_{n+1}\},
  \qquad
  \Ccal_\infty^{(n+1)} \subsetneq \Ccal_\infty^{(n)}.
\]
\end{axiom}

\begin{proposition}[Monotonicity]\label{prop:monotonicity}
The residue is monotone decreasing under expansion:
\[
  \Ccal_\infty^{(1)}
  \;\supseteq\;
  \Ccal_\infty^{(2)}
  \;\supseteq\;
  \cdots
  \;\supseteq\;
  \Ccal_\infty^{(n)}
  \;\supseteq\;
  \cdots
\]
\end{proposition}

\begin{proof}
Immediate: adding a representation adds a constraint, which cannot
enlarge the residue.
\end{proof}

\subsection{The Diagnostic Role of the Residue}

\begin{proposition}[Residue as diagnostic]\label{prop:residue-diagnostic}
The residue $\Ccal_\infty^{(n)}$ identifies the \emph{type} of
information missing at level $n$:
\begin{enumerate}[label=(\roman*)]
  \item if $\Ccal_\infty^{(n)}$ is a single equivalence class, the
        methodology is closed;
  \item if $\Ccal_\infty^{(n)}$ is a $\simadm$-invariant family with
        known parametrisation, the missing information is of
        Type~II (add the parametrising representation);
  \item if $\Ccal_\infty^{(n)}$ is not classifiable within the
        current theory, the missing information is of Type~III
        (a new class of representations is required).
\end{enumerate}
\end{proposition}

\begin{proof}
By inspection of the possibilities.
\end{proof}

\subsection{The Dynamic Diagram}

\begin{equation}
\begin{tikzcd}[row sep=large, column sep=large]
  \Xcal \arrow[r, "\Phimap_{\Rfrak_0}"] \arrow[dr, "\Phimap_{\Rfrak_1}"']
    & \Rcal_0 \arrow[d, "\text{residue}"] \\
  & \Ccal_\infty^{(0)}
    \arrow[d, "\text{expansion}"] \\
  & \Rfrak_1 = \Rfrak_0 \cup \{R_{\rm new}\}
    \arrow[d, "\text{iterate}"] \\
  & \Ccal_\infty^{(1)} \subsetneq \Ccal_\infty^{(0)}
\end{tikzcd}
\end{equation}

\begin{remark}
The dynamic account of representation expansion removes the naive
assumption that the layer set $\Lam$ is fixed in advance. The layer
set is \emph{discovered} through the constraining procedure. This
is a strengthening of the framework paper, in which $\Lam$ was
introduced axiomatically.
\end{remark}

% ------------------------------------------------------------
\section{The Open Problem of Coherent Indistinguishability}
\label{sec:open}
% ------------------------------------------------------------

\subsection{Formulation}

\begin{problem}[Coherent indistinguishability]\label{prob:indist}
Does there exist a class $\Xcal$ of mathematical structures and a
coherent representation family $\Rfrak_{\rm coh}$ such that
\[
  \exists X, Y \in \Xcal : \quad
  X \not\simadm Y, \quad
  \Phimap_{\rm coh}(X) = \Phimap_{\rm coh}(Y),
\]
where $\Phimap_{\rm coh}$ is the coherent representation map of
Definition~\ref{def:coherent-map}?
\end{problem}

\begin{remark}
The problem is not whether a \emph{single} representation fails to
be injective. It asks whether the \emph{entire coherent system} can
fail, i.e.\ whether there exist two genuinely different structures
that every coherently compatible representation family
identifies.
\end{remark}

\subsection{Necessary Conditions for a Counterexample}

\begin{proposition}[Necessary conditions]\label{prop:necessary}
If $X \not\simadm Y$ is a solution to
Problem~\ref{prob:indist}, then:
\begin{enumerate}[label=(\roman*)]
  \item $X$ and $Y$ have the same spectrum under every natural
        Laplacian-type operator;
  \item $X$ and $Y$ are homeomorphic, and in fact homotopy
        equivalent;
  \item $X$ and $Y$ have isomorphic algebraic invariants
        (cohomology rings, $K$-theory, etc.);
  \item the nonlinear structure of any natural evolution equation
        on $X$ and $Y$ agrees to all orders accessible by
        observation;
  \item the numerical behaviour of any reasonable discretisation of
        a natural equation on $X$ and $Y$ agrees asymptotically.
\end{enumerate}
\end{proposition}

\begin{proof}
Each condition is a component of the coherence tuple. If any one
were violated, the corresponding representation would distinguish
$X$ from $Y$, contradicting indistinguishability.
\end{proof}

\begin{remark}
The necessary conditions are extremely restrictive. No known
example satisfies all of them. This suggests, but does not prove,
that Problem~\ref{prob:indist} may have a negative answer for
sufficiently rich $\Xcal$.
\end{remark}

\subsection{Relation to Established Non-Identifiability Results}

\begin{example}[Isospectrality]
Gordon--Webb--Wolpert give Riemannian manifolds $M \not\cong N$
with $\spec(\Delta_M) = \spec(\Delta_N)$. But $M$ and $N$ are
distinguished by their topological layer if their fundamental
groups differ (which in the GWW example they do not, since the
manifolds are homeomorphic), or by their geometric layer if their
metrics differ in curvature. Isospectrality is a Type~II rupture.
\end{example}

\begin{example}[Same linearisation, different nonlinearity]
The equations $u_t = u_{xx} - u^3$ and $u_t = u_{xx} - 2u^3$ have
the same Fr\'echet linearisation at $u = 0$ but are distinguished
by the variational layer. This is also a Type~II rupture.
\end{example}

\begin{example}[Same jets at a point, different global functions]
The functions
\[
  f(x) = \begin{cases} e^{-1/x^2} & x > 0 \\ 0 & x \leq 0
  \end{cases}, \qquad g \equiv 0
\]
have all derivatives equal at $x = 0$ but are globally distinct.
The jet map is not injective; the global analytic layer
distinguishes them. Again Type~II.
\end{example}

\begin{remark}
All known non-identifiability results are Type~II. Problem~\ref{prob:indist}
asks for a Type~III example. No such example is known.
\end{remark}

\subsection{A Conjecture}

\begin{conjecture}[Injectivity of coherent representation]\label{conj:injectivity}
For the class $\Xcal$ of smooth compact Riemannian manifolds
equipped with a natural Laplace-type operator, and for the coherent
representation family
\[
  \Rfrak_{\rm coh}
  =
  \{\text{spectral}, \text{geometric}, \text{topological},
    \text{algebraic}, \text{analytic}\},
\]
the map $\Phimap_{\rm coh}$ is injective modulo isometry.
\end{conjecture}

\begin{remark}
Conjecture~\ref{conj:injectivity} is plausible because the
geometric layer alone is already injective modulo isometry (the
metric determines the manifold). The content of the conjecture is
that the constraining can be \emph{achieved} by the coherent
representation family, not merely by the trivial representation
$g \mapsto g$.
\end{remark}

% ------------------------------------------------------------
\section{Epistemological Implications}
\label{sec:epistemology}
% ------------------------------------------------------------

\subsection{The Two Roles of a Representation}

Each representation $R_i$ plays two distinct roles in the
methodology:
\begin{enumerate}[label=(\roman*)]
  \item a \emph{descriptive} role: it records a piece of structure
        of $X$;
  \item a \emph{constraining} role: it eliminates candidate
        structures incompatible with $r_i$.
\end{enumerate}

The constraining role is the operational content of the
methodology. A representation that is descriptive but not
constraining (e.g.\ a redundant invariant) does not contribute to
separation.

\begin{proposition}[Non-redundancy]\label{prop:non-redundancy}
A representation $R_i$ is \emph{non-redundant} with respect to
$\Rfrak \setminus \{R_i\}$ if and only if
\[
  \Ccal_\infty^{(\Rfrak)} \subsetneq
  \Ccal_\infty^{(\Rfrak \setminus \{R_i\})}.
\]
\end{proposition}

\begin{proof}
Immediate: removing $R_i$ leaves a strictly larger residue.
\end{proof}

\subsection{The Epistemological Limit}

\begin{proposition}[Non-constructive visibility]\label{prop:visibility}
A structure $X \not\simadm Y$ that satisfies
$\Phimap_{\rm coh}(X) = \Phimap_{\rm coh}(Y)$ for a given coherent
family $\Rfrak_{\rm coh}$ is \emph{invisible} to that family. It may
exist without being detectable by any representation in
$\Rfrak_{\rm coh}$.
\end{proposition}

\begin{proof}
By definition of indistinguishability.
\end{proof}

\begin{remark}
This is the precise formal content of the intuition that some
mathematical structures may exist without being epistemologically
accessible until a new class of representations is introduced. The
framework is not committed to the claim that such structures
exist; it merely provides the language in which the question can be
asked.
\end{remark}

\subsection{The Dynamic of Discovery}

The dynamic account of \S\ref{sec:expansion} suggests a
general picture of mathematical discovery:
\[
  \text{problem} \to \text{known representations}
  \to \text{residue} \to \text{new representation}
  \to \text{reduced residue} \to \cdots
\]
The residue is the operational pressure that drives discovery. This
is consistent with historical cases (e.g.\ the development of
topology, group theory, or spectral theory from problems in which
these structures were implicit but not yet formalised).

\subsection{Closing Remark}

The methodology developed in \cite{paperI} and extended here is
best described as a \emph{structural constraining procedure with
dynamic expansion}. Its central criterion is not the existence of
any individual representation, but the \emph{separation} of the
representation family modulo admissible equivalence. Its central
open problem is whether separation can always be achieved, or
whether there exist coherently indistinguishable but
non-isomorphic structures.

The question, in one line:

\begin{quote}
\emph{Can the structure of a mathematical object always be
recovered from its coherent representations, or does mathematics
admit coherently invisible pairs?}
\end{quote}

% ------------------------------------------------------------
\appendix
% ------------------------------------------------------------

\section{Notation Summary}
\label{app:notation}

\begin{center}
\begin{tabular}{c l}
  \toprule
  Symbol & Meaning \\
  \midrule
  $\Xcal$ & Class of mathematical structures \\
  $\Rcal_i$ & Codomain of representation $R_i$ \\
  $\Rfrak$ & Family of representations \\
  $\Phimap$ & Representation map $\Xcal \to \prod_i \Rcal_i$ \\
  $\simadm$ & Admissible equivalence on $\Xcal$ \\
  $\Rfrak_{\rm coh}$ & Coherence subspace of $\prod_i \Rcal_i$ \\
  $\Ccal_\infty^{(n)}$ & Residue at level $n$ \\
  Type I/II/III & Rupture classification \\
  \bottomrule
\end{tabular}
\end{center}

\section{Bibliography}
\label{app:biblio}

\begin{thebibliography}{99}

\bibitem{paperI}
[Companion paper],
\emph{Structural Constraining and Spectral Emergence: A
Multirepresentational Methodology for Nonlinear Evolution
Equations}, 2024.

\bibitem{framework}
[Framework paper],
\emph{A Multirepresentational Construction--Reconstruction
Methodology for Mathematical Structures}, Parts I--V, 2024.

\bibitem{gordon}
C.~Gordon, D.~Webb, S.~Wolpert,
Isospectral plane domains and surfaces via Riemannian orbifolds,
\emph{Invent.~Math.}~110 (1992), 1--22.

\bibitem{kac}
M.~Kac,
Can one hear the shape of a drum?
\emph{Amer.~Math.~Monthly}~73 (1966), 1--23.

\bibitem{engl}
H.~W.~Engl, M.~Hanke, A.~Neubauer,
\emph{Regularization of Inverse Problems}, Kluwer, 1996.

\bibitem{reed-simon-1}
M.~Reed, B.~Simon,
\emph{Methods of Modern Mathematical Physics, Vol.~I},
Academic Press, 1980.

\end{thebibliography}

\end{document}